% ===================================================================
%  LENTELĖ Nr. 14 — Gatvė. --- Prekybininkai.
%  Traduction 2026 d'apres texto/io, controlee sur texto/fr, texto/ru
%  et texto/cs. Consigne : voir texto/lt/10-lentele-01.tex.
%
%  SOIXANTE METIERS EN SEPT PAGES, ET LE TABLEAU LE PLUS PEUPLE DU
%  LIVRET APRES LE 6. Il donne a la colonne lituanienne sa
%  demonstration la plus nette, parce que la langue nomme un metier
%  AVEC QUATRE SUFFIXES ET RIEN D'AUTRE :
%
%    -ininkas, tire de la CHOSE : « laikrodininkas » l'horloger de
%      « laikrodis », « skėtininkas » le marchand de parapluies de
%      « skėtis », « ginklininkas » l'armurier de « ginklas »,
%      « spaustuvininkas » l'imprimeur de « spaustuvė » ;
%    -ėjas / -tojas, tire du VERBE : « kirpėjas » celui qui coupe,
%      « siuvėjas » celui qui coud, « pardavėjas » celui qui vend,
%      « leidėjas » celui qui edite, « raižytojas » celui qui grave ;
%    -ys, sur un COMPOSE : « batsiuvys » qui coud les bottes,
%      « knygrišys » qui relie les livres, « kaminkrėtys » qui ramone
%      les cheminees, « auksakalys » qui forge l'or ;
%    -vys / -ius, sur le METAL : « kalvis » le forgeron, et
%      « šaltkalvis » le serrurier — mot a mot le forgeron-a-froid.
%
%  DE SORTE QU'ON PEUT FABRIQUER LE NOM D'UN METIER SANS L'AVOIR
%  APPRIS. C'est le prolongement direct du tableau 6, ou l'ustensile
%  se nommait par son action ; ici c'est l'HOMME qui se nomme par la
%  sienne.
%
%  ET LES EXCEPTIONS TOMBENT TOUTES DU MEME COTE. « fotografas »,
%  « optikas », « dentistas », « advokatas », « notaras »,
%  « konsulas », « juvelyras », « stenografas », « generolas »,
%  « majoras », « kapitonas », « leitenantas » restent latins. LE
%  METIER QUI S'APPREND DANS UN ATELIER RECOIT UN NOM LITUANIEN, LA
%  PROFESSION QUI S'APPREND SUR UN DIPLOME GARDE LE SIEN. C'est le
%  critere du lieu d'apprentissage, invente pour un mot irlandais au
%  tableau 10 de la onzieme colonne, applique ici a soixante d'un
%  coup.
%
%  ET LA MEME LIGNE COUPE LES GRADES. « pulkininkas » le colonel est
%  lituanien, tire de « pulkas » le regiment, et « puskarininkiai »
%  les sous-officiers en est un compose ; « majoras »,
%  « kapitonas », « leitenantas », « generolas » viennent du dehors.
%  Le regiment est vieux, l'etat-major est neuf.
%
%  QUATRE GENITIFS TENUS PAR UNE RELATIVE OU UNE APPOSITION : le
%  bureau (32) du cousin (34), le jet d'eau (49) de la fontaine (48),
%  l'imperiale (59) de l'omnibus (58).
% ===================================================================

%%K t14-apar-1 apar
\VUsaut{4.0mm}
\VUtitre{15.0pt}{60}{LENTELĖ Nr. 14}
\VUsaut{3.0mm}

%%K t14-tit-1 sub
\VUpk{11.4pt}{40}{Gatvė. --- Prekybininkai.}
\VUsaut{3.0mm}

%%K t14-01-1 p
1. --- Mano draugas Jonas, kuris gyvena kaime, atvyko praleisti trijų dienų mano
šeimoje. \VUgras{Iki} tol jis neturėjo progos pamatyti didelio miesto; todėl visas
dienas praleidžiame vaikščiodami, o aš aiškinu jam tai, ko jis nepažįsta.

%%K t14-02-1 p
2. --- Pirmiausia nuėjome aplankyti \VUgras{observatorijos} \textsuperscript{(1)},
kuri jį labai sudomino. Grįždami per \VUgras{gatvę} \textsuperscript{(3)}, kurioje yra
\VUgras{Turkiška pirtis} \textsuperscript{(2)}, sutikome \VUgras{laidotuves}
\textsuperscript{(4)}. Prieš \VUgras{laidotuvių eiseną} ėjo \VUgras{dvasininkai,}
einantys pirma \VUgras{katafalko,} kuriame buvo padėtas \VUgras{karstas.} Daug draugų
lydėjo \VUgras{gedinčią} šeimą.

%%K t14-02-2 p
Palydai praėjus, perėjome gatvę prieš didelę \VUgras{alinę} \textsuperscript{(5)},
kurioje daug \VUgras{lankytojų} gėrė \VUgras{alų} terasoje, prisidengę stikline
\VUgras{markize} \textsuperscript{(6)}.

%%K t14-03-1 p
3. --- Jonas buvo apstulbintas \VUgras{herbo skydo,} pakabinto tarp \VUgras{likerių
pardavėjo} \textsuperscript{(7)} krautuvės ir \VUgras{muzikos pardavėjo}
\textsuperscript{(10)} krautuvės. Jis paklausė manęs, ką jis reiškia; atsakiau, kad jis
žymi anglų \VUgras{konsulatą} \textsuperscript{(8)}, ir parodžiau jam \VUgras{konsulą}
\textsuperscript{(9)}, kuris yra balkone.

%%K t14-tit-2 sub
\VUpk{10.2pt}{40}{Krautuvių apžiūra.}
\VUsaut{3.0mm}

%%K t14-04-1 p
4. --- Savaime suprantama, sustojome prieš \VUgras{muzikos instrumentų,}
\VUgras{fisharmonijų,} \VUgras{vargonų} ir \VUgras{pianinų} vitriną; žiūrėjome į
iliustruotus \VUgras{partitūrų} ir pianino \VUgras{natų} viršelius ir gėrėjomės
\VUgras{lyra} iš auksuoto medžio, kuri naudojama kaip iškaba.

%%K t14-05-1 p
5. --- Dulkių debesys, keliami \VUgras{šlavėjų} \textsuperscript{(11)} ir
\VUgras{šlavimo vežimo} \textsuperscript{(12)}, privertė mus greitai praeiti pro gražius
\VUgras{baldų pardavėjo} \textsuperscript{(13)} \VUgras{baldus} ir pro
\VUgras{skardininko} \textsuperscript{(14)} \VUgras{krosnelių} ir \VUgras{lempų}
vitriną.

%%K t14-06-1 p
6. --- Beje, šioje vietoje turėjome palikti šaligatvį, nes jis buvo užgriozdintas
ryšuliais, kuriuos ėmė iš \VUgras{pervežimo vežimo} \textsuperscript{(16)}, kad sudėtų
juos į \VUgras{sandėlį} \textsuperscript{(15)}, kurį ką tik buvo išsinuomoję.

%%K t14-06-2 p
Tai sukliudė mums pamatyti gražius \VUgras{ekipažų dirbėjo} \textsuperscript{(17)}
ekipažus, bet nesukliudė sustoti pažiūrėti į \VUgras{pakuotoją}
\textsuperscript{(19)}, darantį dėžę savo kaimynui, \VUgras{dviračių ir automobilių
pardavėjui} \textsuperscript{(18)}. Paskui, kadangi jau buvo kiek vėlu, grįžome namo.

%%K t14-tit-3 sub
\VUpk{10.2pt}{40}{Pasivaikščiojimas po miestą.}
\VUsaut{3.0mm}

%%K t14-07-1 p
7. --- Kitą dieną mano motina pasiūlė Jonui nusifotografuoti, paskui pavedė man
palydėti jį pas \VUgras{fotografą} \textsuperscript{(20)}. Po pusryčių nuėjome į
menininko \VUgras{dirbtuvę,} kur turėjome gana ilgai laukti. Kad užsiimtume, žiūrėjome
į \VUgras{portretus} \textsuperscript{(22)}, kabančius ant sienos.

%%K t14-07-2 p
Kai atėjo mūsų eilė, darbas truko neilgai, ir vos mano draugas baigė pozuoti prieš
\VUgras{fotoaparatą} \textsuperscript{(21)}, nusileidome.

%%K t14-08-1 p
8. --- Antrajame aukšte pro \VUgras{kirpėjo} \textsuperscript{(23)} duris, kurios
atvertos, matėme jo pameistrį, kertantį kliento plaukus.

%%K t14-08-2 p
Atėję į gatvę, praėjome pro \VUgras{tabako krautuvę} \textsuperscript{(24)}. Joną taip
palinksmino raudonas \VUgras{žibintas} \textsuperscript{(25)} ir \VUgras{stora pypkė,}
kuri naudojama kaip \VUgras{iškaba} \textsuperscript{(26)}, kad jis susidūrė su dviem
senais ponais, kurie kalbėjosi \VUgras{uostydami tabaką} \textsuperscript{(27)}.

%%K t14-tit-4 sub
\VUpk{10.2pt}{40}{Konditerija.}
\VUsaut{3.0mm}

%%K t14-09-1 p
9. --- Visų šalių \VUgras{banknotai} ir \VUgras{monetos,} išdėti \VUgras{keitėjo}
\textsuperscript{(29)} \VUgras{vitrinoje,} akimirkai patraukė mūsų dėmesį, bet kadangi
buvo užkandžių valanda, įėjome pas \VUgras{konditerį} \textsuperscript{(31)} ilgiau
nestodami.

%%K t14-10-1 p
10. --- Matydami visus \VUgras{skanėstus,} kurių buvo krautuvėje ---
\VUgras{pyragaičius, meduolius, sausainius} ir visokius \VUgras{saldainius} ---,
sunkiai galėjome išsirinkti. Galiausiai Jonas išsirinko \VUgras{kreminį pyragaitį,} o
aš paėmiau tik mažą \VUgras{vyšnių pyragėlį,} nes saldumynai \VUgras{gadina man
dantis,} ir aš visai nenoriu per dažnai lankytis pas \VUgras{odontologą}
\textsuperscript{(30)}.

%%K t14-tit-5 sub
\VUpk{10.2pt}{40}{Rašymas mašinėle.}
\VUsaut{3.0mm}

%%K t14-11-1 p
11. --- Norėdami parodyti savo draugui \VUgras{rašomąją mašinėlę,} apie kurią jis buvo
girdėjęs, bet kurios nepažino, įėjome į \VUgras{kontorą} \textsuperscript{(32)}, kuri
priklauso mano \VUgras{pusbroliui} \textsuperscript{(34)}. Radome jį diktuojantį savo
laiškus \VUgras{sekretoriui} \textsuperscript{(35)}, kuris yra \VUgras{stenografas.} Vos
laiškas baigtas, \VUgras{mašininkė} \textsuperscript{(33)} perrašydavo jį savo mašinėle.
Nenorėdami trukdyti savo giminaičio užsiėmimų, pabuvome pas jį tik kelias akimirkas.

%%K t14-12-1 p
12. --- Po mano pusbrolio kontora gyvena didelis \VUgras{laikrodininkas-juvelyras}
\textsuperscript{(37)}, kuris dar yra ir auksakalys. Jo puiki krautuvė turi daug
\VUgras{laikrodžių, švytuoklinių laikrodžių, kišeninių laikrodžių,}
\VUgras{grandinėlių} ir brangių \VUgras{brangenybių,} pavyzdžiui, \VUgras{perlų
vėrinių,} auksinių \VUgras{apyrankių, auskarų} ir \VUgras{žiedų,} papuoštų
\VUgras{brangakmeniais} ir \VUgras{deimantais}. Viena vitrina ypač skirta
\VUgras{auksiniams dirbiniams} ir turi svarbų sidabrinių \VUgras{stalo įrankių}
rinkinį.

%%K t14-13-1 p
13. --- Sukdamas už gatvės kampo, pastebėjau, kad buvo pakeista \VUgras{lentelė}
\textsuperscript{(36)}, pakabinta prie namo \VUgras{numerio.} Jau garsiai ketinau
pasakyti savo pastabą, kai pasigirdo džiaugsmingi \VUgras{karinės} muzikos garsai.
Nubėgome pulkui priešais, nepagalvoję, kad \VUgras{jis} eis pro
\VUgras{skrybėlininko} \textsuperscript{(39)} ir \VUgras{pirštininko}
\textsuperscript{(40)} krautuves ir kad būtume buvę lygiai taip pat gerai įsitaisę
matyti jo paradą, likę prieš \VUgras{optiko} \textsuperscript{(38)} vitriną, pilną
\VUgras{pensnė, akinių} ir \VUgras{žiūronėlių}.

%%K t14-tit-6 sub
\VUpk{10.2pt}{40}{Žygiavimo paradas.}
\VUsaut{3.0mm}

%%K t14-14-1 p
14. --- Prieš \VUgras{pulką} \textsuperscript{(41)} žygiavo \VUgras{būgnininkų vadas}
\textsuperscript{(42)}, lydimas \VUgras{būgnininkų} \textsuperscript{(43)},
\VUgras{trimitininkų} \textsuperscript{(44)} ir viso \VUgras{orkestro.}
\VUgras{Generolas} \textsuperscript{(45)}, kuris buvo aikštėje su savo adjutantu, buvo
sustojęs prie \VUgras{vandens čiurkšlės} \textsuperscript{(49)}, kuri trykšta iš
\VUgras{fontano} \textsuperscript{(48)}, kad stebėtų \VUgras{paradą.}

%%K t14-14-2 p
\VUgras{Štabo karininkai} \textsuperscript{(46)}, \VUgras{pulkininkas} ir
\VUgras{pulkininkas leitenantas} pasveikino jį špaga. \VUgras{Majorai} buvo prieš savo
\VUgras{batalionus} \textsuperscript{(47)}, ir kiekvienas \VUgras{kapitonas} vadovavo
savo \VUgras{kuopai,} o \VUgras{leitenantai,} \VUgras{jaunesnieji leitenantai} ir
\VUgras{puskarininkiai} užėmė savo įprastą vietą prie savo vyrų.

%%K t14-15-1 p
15. --- Daug praeivių susibūrė \VUgras{sankryžoje} \textsuperscript{(50)};
prekybininkai --- \VUgras{skėtininkas} \textsuperscript{(51)}, \VUgras{audinių
dažytojas} \textsuperscript{(53)}, \VUgras{ginklininkas} \textsuperscript{(54)} ---
išėjo pamatyti \VUgras{kareivių.} Visos priemonės --- \VUgras{fiakrai}
\textsuperscript{(52)}, \VUgras{ponų ekipažai} \textsuperscript{(57)} ir net
\VUgras{laistymo statinė} \textsuperscript{(56)} --- sustojo palei šaligatvį.

%%K t14-tit-7 sub
\VUpk{10.2pt}{40}{Scenos gatvėje.}
\VUsaut{3.0mm}

%%K t14-15-2 p
\VUgras{Prekybos keliauninkas} \textsuperscript{(55)}, nešantis rankoje savo
\VUgras{pavyzdžių dėžutes,} greitai perėjo gatvę, o žmonės, sėdintys ant
\VUgras{imperialo} \textsuperscript{(59)}, \VUgras{omnibuso} \textsuperscript{(58)}
viršutinio aukšto, atsisuko pasigėrėti reginiu. Vargšas \VUgras{klajojantis
dainininkas} \textsuperscript{(60)} su savo \VUgras{ryla} \textsuperscript{(61)} turėjo
nutraukti savo \VUgras{dainą,} nes būgnų triukšmas užgožė jo balsą.

%%K t14-16-1 p
16. --- Kai pulkas atėjo prieš \VUgras{audinių krautuvę} \textsuperscript{(64)},
\VUgras{siuvėjos} \textsuperscript{(63)} ir \VUgras{sukirpėjai} \textsuperscript{(62)}
paliko savo \VUgras{siuvimo mašinas} ir savo darbą, kad pažiūrėtų pro langą.
\VUgras{Prekybininkas} \textsuperscript{(65)}, kuris ką tik buvo iškabinęs naujus
\VUgras{kostiumus} \textsuperscript{(66)} savo \VUgras{vitrinoje,} turėjo liepti sunešti
vidun \VUgras{gelumbes} \textsuperscript{(67)} ir \VUgras{audinius,} išdėtus ant
šaligatvio prie \VUgras{kanalizacijos angos} \textsuperscript{(68)}, bijodamas, kad
minia juos nuvers; net senas \VUgras{prekiautojas iš rankų} \textsuperscript{(69)}, iš
kurio durininkė derėjosi dėl kojinių, buvo priverstas įeiti į koridorių, kad baigtų
savo pardavimą.

%%K t14-16-2 p
Palydėjome pulką iki kareivinių \VUgras{ir,} labai patenkinti savo diena, grįžome
vakarienės valandą.

%%K t14-tit-8 sub
\VUpk{10.2pt}{40}{Įvairūs verslai ir profesijos.}
\VUsaut{3.0mm}

%%K t14-17-1 p
17. --- Kitą dieną mano tėvas pasiūlė mums aplankyti tenykščio laikraščio spaustuvę.
Su entuziazmu sutikome.

%%K t14-17-2 p
\VUgras{Spaustuvininkas} yra ir \VUgras{knygininkas} \textsuperscript{(70)}
\VUgras{(= knygų pardavėjas),} \VUgras{leidėjas,} \VUgras{popieriaus pardavėjas} ir
\VUgras{knygrišys}. Antrajame aukšte jis įrengė laikraščio \VUgras{redakciją}
\textsuperscript{(71)}, taip pat \VUgras{raižyklą,} kurioje dirba
\VUgras{litografai} \textsuperscript{(72)} ir \VUgras{vario raižytojai}
\textsuperscript{(73)}, galiausiai \VUgras{rinkyklą,} kurioje yra \VUgras{spaustuvės
darbininkai} \textsuperscript{(74)}. Tikroji \VUgras{spausdinimo salė}
\textsuperscript{(75)}, \VUgras{spaudos presai} ir įvairios mašinos yra pirmajame
aukšte.

%%K t14-tit-9 sub
\VUpk{10.2pt}{40}{Jonas labai daug klausinėja.}
\VUsaut{3.0mm}

%%K t14-18-1 p
18. --- Išeidamas iš spaustuvės, Jonas labai nustebęs sustojo prieš mažą stiklinę
dėžutę, pastatytą ant stulpo prie \VUgras{galanterijos krautuvės}
\textsuperscript{(77)}, ir paklausė, kam ji naudojama. Mano tėvas paaiškino jam, kad
tai \VUgras{gaisro pranešiklis} \textsuperscript{(76)}. Beje, viskas mieste mano draugui
buvo nuostabu --- nuo paprasto \VUgras{akmeninio fontano} \textsuperscript{(78)} ir
lengvo \VUgras{krovininio triračio} \textsuperscript{(80)}, kurį vairuoja
\VUgras{pasiuntinukas} \textsuperscript{(79)} su livrėja, iki \VUgras{pašto vežimo}
\textsuperscript{(81)}.

%%K t14-19-1 p
19. --- Kol mano tėvas kelioms akimirkoms paliko mus pasikalbėti su \VUgras{mokesčių
rinkėju} \textsuperscript{(83)}, kuris buvo sustojęs pasišnekėti su \VUgras{advokatu}
\textsuperscript{(82)} ir \VUgras{notaru} \textsuperscript{(84)}, mes linksminomės
sekdami akimis judėjimą gatvėje. Pro mus praėjo patys įvairiausi žmonės:
\VUgras{konditerio pameistrys} \textsuperscript{(85)}, nešantis krepšyje puikų
\VUgras{paštetą,} atėjo paskui \VUgras{drabužių pardavėją} \textsuperscript{(86)};
\VUgras{žibintų uždegėjas} \textsuperscript{(87)} kalbėjosi su \VUgras{dujininku}
\textsuperscript{(88)}; \VUgras{vienuolė} \textsuperscript{(89)}, laikanti už rankos
našlaitę, grįžo į savo vienuolyną.

%%K t14-tit-10 sub
\VUpk{10.2pt}{40}{Grįžtant namo.}
\VUsaut{3.0mm}

%%K t14-20-1 p
20. --- Tą akimirką, kai mano tėvas mus pasivijo, Jonas kvatojo, matydamas nešvarius
veidus ir keistus apdarus dviejų visai suodinų \VUgras{kaminkrėčių}
\textsuperscript{(91)}.

%%K t14-21-1 p
21. --- Norėjau nupirkti iš \VUgras{gėlininkės} \textsuperscript{(92)} augalą savo
motinai, bet mano tėvas parodė man, kad jis būtų labai sunkus nešti ir kad būtų
geriau įsigyti \VUgras{apelsinų}. Priėjome prie \VUgras{pardavėjos}
\textsuperscript{(94)}, ir kai \VUgras{auklė} \textsuperscript{(93)}, kuri rinkosi savo
auklėtiniui, baigė pirkti, liepėme mus aptarnauti.

%%K t14-22-1 p
22. --- Grįždami namo sutikome kelis pažįstamus: seną mano šeimos draugę, kuri rėmėsi
į savo \VUgras{palydovę} \textsuperscript{(95)}, tiltų ir kelių \VUgras{inžinierių}
\textsuperscript{(96)} ir vieną savo pusseserę su jos \VUgras{vaikų aukle}
\textsuperscript{(98)}.

%%K t14-22-2 p
Paskui sutikome mano motinos \VUgras{skrybėlininkę} \textsuperscript{(97)} ir du
miestui svetimus \VUgras{vienuolius} \textsuperscript{(99)}, kurie priėjo prie mūsų
paklausti mano tėvo kelio į stotį.

\VUsaut{6.0mm}
\VUfilet{22mm}
